\chapter{Anexo B: pipeline de datos}

Este anexo recoge la documentacion de \path{doc/docs/data/}: vision general del ETL, lectura raw, union de mercado y contratos, separacion por producto, asociacion del subyacente y calculo de volatilidad implicita.

\section{Vision general del ETL}

El ETL transforma contratos, operaciones y tipos en \path{VOLATILITY_DB}. Cada paso consume bases previas, aplica validaciones y materializa una salida. Las entradas conceptuales son contratos \path{CCONTRACTS_C2}, trades \path{TGENTRADES} y series EONIA/euro short-term rate. Las salidas principales son contratos IBEX limpios, trades IBEX, opciones, futuros, relacion opcion-subyacente, opciones con subyacente y volatilidad implicita.

\section{Read Raw Step}

La lectura raw normaliza cabeceras y tipos, filtra contratos del universo IBEX y construye la base de tipos. El documento \path{raw-read.md} detalla entradas, normalizacion, filtrado y validaciones. La decision clave es que los errores de formato se corrigen o descartan al principio, antes de que se mezclen con calculos financieros.

\section{Merge Raw Step}

\path{merge-raw.md} documenta la union por contrato y fecha. Tambien describe imputacion de vencimientos, imputacion de strikes y calculo de tiempo hasta vencimiento. La coherencia entre codigo de contrato, strike y vencimiento se verifica con la configuracion de codigos.

\section{Product Split Step}

\path{product-split.md} separa opciones y futuros. Las opciones se renombran con \path{OptionContractCode}; los futuros, con \path{FutureContractCode}. La relacion opcion-subyacente se construye a partir de madurez y convenciones de contrato. Esta relacion permite buscar despues el futuro temporalmente valido para cada opcion.

\section{Underlying Step}

\path{underlying.md} describe la union temporal as-of entre opcion y futuro. El criterio es usar el futuro mas cercano anterior o simultaneo, evitando informacion futura. La salida incluye \path{UnderlyingPrice}, \path{UnderlyingExecDatetime} y \path{UnderlyingLagMinutes}. La base intermedia contiene 428.448 filas antes del filtrado final de volatilidad.

\section{Volatility Step}

\path{volatility.md} documenta limpieza de operaciones, validaciones, tipo compuesto, formula Black-76, solver de volatilidad implicita, filtro por tipo de operacion y salida final. La base final contiene 312.615 filas, 151.424 calls y 161.191 puts. El filtro de lag maximo de 180 minutos y la resolucion numerica eliminan observaciones no aptas.

\section{Columnas finales}

\path{VOLATILITY_DB} contiene fecha, codigo de opcion, tipo call/put, mercado, identificador de trade, precio de opcion, cantidad, tipo de trade, strike, madurez, fecha-hora de ejecucion, tiempo a vencimiento, tipo, codigo de futuro, precio del subyacente, fecha-hora del subyacente, lag y volatilidad implicita.
